第一部分：基于线图理论的非加性最短路径问题建模与 CWSP 算法框架设计。深入剖析非加性约束条件引起算法回溯及解空间发散的核心根源。利用图论空间映射原理，详细建立由底向上的原图 $G(V,E)$ 至线图 $L(G)$ 的标准化拓扑转换流程，并提出条件带权最短路径（CWSP）算法。明确论证了原网络中的边界向线图节点过渡的邻接点多重语义机制，构建出适应动态转角约束的惩罚权值分配模型，从理论上成功规避由分支夹角引起的连贯通行障碍。

第二部分：跨越算力瓶颈的 CWSP 算法 GPU 并行化架构优化。针对线图重构带来的 $O(m^2)$ 规模膨胀与实时寻路时延问题，深入挖掘 CUDA 的细粒度并行潜力，对 CWSP 算法进行了底层结构重塑。在拓扑构建期，引入两阶段前缀和与无锁 CSR 写入模式，消除边裂变期的高频内存开辟碎片；在路径寻优期，以 $\Delta$-Stepping 算法替代传统的优先队列，通过桶式并行迭代与原子归约指令，彻底打破管线中的锁竞争限制，实现了从图预处理到波前松弛寻优的全链路无锁高并发执行。

第三部分：在大规模交通与复杂网络环境下的数值实验与加速应用评估。为验证并行化 CWSP 算法的准确性及其在现实极端场景下的工程有效性，本文在真实的微观机场起降网络平台与自合成的大范围致密环境中执行了体系化的基准测试。借助毫秒级时间剖析，详细探讨了拓扑稠密度、跨步步长（$\Delta$）等核心调优参数对加速曲线的敏感机制，最终通过对比数据明确论证了本文提出的异构求解模型在应对真实非加性调度时的绝对主导性能优势。
